\documentclass[12pt,a4paper]{article}

% ── Packages ──────────────────────────────────────────────────────────────────
\usepackage[utf8]{inputenc}
\usepackage[T1]{fontenc}
\usepackage[french]{babel}
\usepackage{geometry}
\usepackage{graphicx}
\usepackage{xcolor}
\usepackage{amsmath,amssymb}
\usepackage{booktabs}
\usepackage{array}
\usepackage{tabularx}
\usepackage{multirow}
\usepackage{longtable}
\usepackage{hyperref}
\usepackage{listings}
\usepackage{fancyhdr}
\usepackage{titlesec}
\usepackage{tcolorbox}
\usepackage{enumitem}
\usepackage{float}
\usepackage{caption}
\usepackage{subcaption}
\usepackage{mdframed}
\usepackage{microtype}
\usepackage{parskip}

% ── Page geometry ─────────────────────────────────────────────────────────────
\geometry{
  top=2.5cm, bottom=2.5cm,
  left=2.8cm, right=2.8cm,
  headheight=14pt
}

% ── Colours ───────────────────────────────────────────────────────────────────
\definecolor{primary}{RGB}{30,90,160}
\definecolor{accent}{RGB}{220,60,30}
\definecolor{lightgray}{RGB}{245,245,245}
\definecolor{darkgray}{RGB}{80,80,80}
\definecolor{codegreen}{RGB}{0,128,0}
\definecolor{codepurple}{RGB}{128,0,128}
\definecolor{codegray}{RGB}{120,120,120}
\definecolor{backcolour}{RGB}{248,248,248}
\definecolor{successgreen}{RGB}{34,139,34}
\definecolor{warnorange}{RGB}{220,100,0}

% ── Section styles ────────────────────────────────────────────────────────────
\titleformat{\section}
  {\Large\bfseries\color{primary}}
  {\thesection}{1em}{}[\titlerule]

\titleformat{\subsection}
  {\large\bfseries\color{primary!80!black}}
  {\thesubsection}{1em}{}

\titleformat{\subsubsection}
  {\normalsize\bfseries\color{darkgray}}
  {\thesubsubsection}{1em}{}

% ── Code listings ─────────────────────────────────────────────────────────────
\lstdefinestyle{pythonstyle}{
  backgroundcolor=\color{backcolour},
  commentstyle=\color{codegreen}\itshape,
  keywordstyle=\color{primary}\bfseries,
  numberstyle=\tiny\color{codegray},
  stringstyle=\color{codepurple},
  basicstyle=\ttfamily\footnotesize,
  breakatwhitespace=false,
  breaklines=true,
  captionpos=b,
  keepspaces=true,
  numbers=left,
  numbersep=5pt,
  showspaces=false,
  showstringspaces=false,
  showtabs=false,
  tabsize=2,
  frame=single,
  rulecolor=\color{primary!30}
}
\lstset{style=pythonstyle}

% ── Fancy boxes ───────────────────────────────────────────────────────────────
\tcbuselibrary{skins,breakable}

\newtcolorbox{resultbox}[1][]{
  enhanced,
  colback=primary!5,
  colframe=primary,
  fonttitle=\bfseries,
  title={#1},
  breakable
}

\newtcolorbox{warnbox}[1][]{
  enhanced,
  colback=warnorange!10,
  colframe=warnorange,
  fonttitle=\bfseries,
  title={#1},
  breakable
}

\newtcolorbox{successbox}[1][]{
  enhanced,
  colback=successgreen!8,
  colframe=successgreen,
  fonttitle=\bfseries,
  title={#1},
  breakable
}

% ── Header / Footer ───────────────────────────────────────────────────────────
\pagestyle{fancy}
\fancyhf{}
\fancyhead[L]{\small\color{darkgray} Analyse Comportementale Clientèle Retail}
\fancyhead[R]{\small\color{darkgray} Machine Learning — GI2S1}
\fancyfoot[C]{\small\thepage}
\renewcommand{\headrulewidth}{0.4pt}
\renewcommand{\footrulewidth}{0pt}

% ── Hyperref ─────────────────────────────────────────────────────────────────
\hypersetup{
  colorlinks=true,
  linkcolor=primary,
  citecolor=primary,
  urlcolor=accent,
  pdftitle={Rapport ML Retail — Analyse Comportementale Clientèle},
  pdfauthor={Module Machine Learning — GI2S1 MolkaChaoui}
}

% ── Misc ──────────────────────────────────────────────────────────────────────
\captionsetup{font=small,labelfont=bf}
\setlength{\parskip}{6pt}

% ==============================================================================
\begin{document}
% ==============================================================================

% ── Title page ────────────────────────────────────────────────────────────────
\begin{titlepage}
  \centering
  \vspace*{1.5cm}

  {\color{primary}\rule{\linewidth}{2pt}}
  \vspace{0.5cm}

  {\Huge\bfseries\color{primary}
    Analyse Comportementale\\[0.3em]
    Clientèle Retail}

  \vspace{0.4cm}
  {\color{primary}\rule{\linewidth}{0.8pt}}

  \vspace{0.8cm}
  {\LARGE\bfseries Rapport de Projet Machine Learning}

  \vspace{1.2cm}

  \begin{tcolorbox}[
    enhanced, width=0.85\linewidth,
    colback=lightgray,
    colframe=primary,
    fontupper=\large
  ]
  \centering
  Système complet de Machine Learning pour la \textbf{rétention client},
  la \textbf{segmentation comportementale} et la \textbf{prévision des dépenses}
  — déployé via une application Flask.
  \end{tcolorbox}

  \vspace{1.5cm}

  

  \vspace{1.8cm}

  \begin{tabular}{ccc}
    \begin{tcolorbox}[width=3.5cm,colback=successgreen!15,colframe=successgreen,halign=center]
      \textbf{AUC-ROC}\\[2pt]
      {\Large\bfseries\color{successgreen} 0.958}\\
      {\small Classification}
    \end{tcolorbox}
    &
    \begin{tcolorbox}[width=3.5cm,colback=primary!15,colframe=primary,halign=center]
      \textbf{Silhouette}\\[2pt]
      {\Large\bfseries\color{primary} 0.228}\\
      {\small Clustering}
    \end{tcolorbox}
    &
    \begin{tcolorbox}[width=3.5cm,colback=warnorange!15,colframe=warnorange,halign=center]
      \textbf{R²}\\[2pt]
      {\Large\bfseries\color{warnorange} 0.668}\\
      {\small Régression}
    \end{tcolorbox}
  \end{tabular}

  \vfill
  {\color{primary}\rule{\linewidth}{1.5pt}}
  \vspace{0.3cm}
  {\small\color{darkgray} Atelier Machine Learning — Analyse Comportementale Clientèle Retail (UK)}
\end{titlepage}

% ── Table of Contents ────────────────────────────────────────────────────────
\tableofcontents
\newpage

% ==============================================================================
\section{Introduction et Vue d'Ensemble}
% ==============================================================================

\subsection{Contexte du Projet}

Ce projet s'inscrit dans le cadre du module \textbf{Machine Learning — GI2S1} et vise à
concevoir un pipeline bout-en-bout d'analyse comportementale pour une entreprise de retail
e-commerce basée au Royaume-Uni. La compréhension fine du comportement client est aujourd'hui
un enjeu stratégique majeur : identifier les clients à risque de départ (\textit{churn}),
segmenter la base client et estimer la valeur monétaire future permettent d'orienter
les actions marketing et d'optimiser le chiffre d'affaires.

Le système développé répond à \textbf{trois questions métier critiques} :

\begin{enumerate}[label=\textbf{Q\arabic*.}]
  \item \textbf{Ce client va-t-il partir ?} — Prédiction du churn via un modèle de classification.
  \item \textbf{À quel segment appartient-il ?} — Segmentation non supervisée par clustering.
  \item \textbf{Combien va-t-il dépenser ?} — Estimation de la valeur client par régression.
\end{enumerate}

\subsection{Résumé des Performances}

\begin{table}[H]
\centering
\caption{Récapitulatif des performances des trois modèles}
\begin{tabular}{llll}
\toprule
\textbf{Question} & \textbf{Modèle} & \textbf{Métrique clé} & \textbf{Valeur} \\
\midrule
Prédiction churn & Random Forest + SMOTE & AUC-ROC & \textbf{0.958} \\
                 &                        & F1-Score & \textbf{0.87} \\
\midrule
Segmentation     & KMeans ($k=4$)         & Silhouette Score & \textbf{0.228} \\
                 &                        & Davies-Bouldin   & 1.30 \\
\midrule
Prévision valeur & Régression Linéaire    & $R^2$ & \textbf{0.668} \\
                 &                        & MAE   & 864 £ \\
\bottomrule
\end{tabular}
\end{table}

\subsection{Architecture Générale du Projet}

Le projet est organisé selon la structure suivante :

\begin{mdframed}[backgroundcolor=lightgray, linecolor=primary, linewidth=1pt]
\begin{verbatim}
projet_ml_retail/
├── data/
│   ├── raw/        # Données brutes (CSV source)
│   ├── processed/  # Données après prétraitement
│   └── train_test/ # Données de train/test splittées
├── src/
│   ├── preprocessing.py   # Pipeline nettoyage (6 étapes)
│   ├── train_model.py     # Entraînement des 3 modèles
│   ├── predict.py         # Inférence sur nouvelles données
│   └── utils.py           # Fonctions utilitaires
├── app/
│   ├── app.py             # Application Flask
│   └── templates/         # Interface HTML/Bootstrap
├── notebooks/             # Analyses exploratoires (5 notebooks)
└── requirements.txt
\end{verbatim}
\end{mdframed}

% ==============================================================================
\section{Données et Analyse Exploratoire}
% ==============================================================================

\subsection{Description du Dataset}

Le jeu de données comprend \textbf{4\,372 clients} issus d'une plateforme e-commerce britannique,
décrits par \textbf{52 features initiales} couvrant plusieurs dimensions comportementales :

\begin{itemize}
  \item \textbf{Données démographiques} : âge, genre, pays, région.
  \item \textbf{Comportement d'achat (RFM)} : récence, fréquence, montant total (\textit{MonetaryTotal}).
  \item \textbf{Préférences} : saison favorite, mois préféré, heure préférée, mode de paiement.
  \item \textbf{Service client} : tickets support, score de satisfaction (1--5).
  \item \textbf{Features dérivées} : panier moyen (\textit{AvgBasketValue}), ratio de retour,
        délai moyen entre achats, ratio d'achat le week-end.
  \item \textbf{Variable cible} : \textit{Churn} — variable binaire (0 = fidèle, 1 = parti).
\end{itemize}

\subsection{Distribution de la Variable Cible}

La distribution du churn présente un \textbf{déséquilibre modéré} :

\begin{table}[H]
\centering
\caption{Distribution de la variable cible Churn}
\begin{tabular}{lcc}
\toprule
\textbf{Classe} & \textbf{Effectif} & \textbf{Proportion} \\
\midrule
Churn = 0 (fidèle) & 2\,918 & 66.7\% \\
Churn = 1 (parti)  & 1\,454 & 33.3\% \\
\midrule
\textbf{Total}     & \textbf{4\,372} & \textbf{100\%} \\
\bottomrule
\end{tabular}
\end{table}

Ce déséquilibre (ratio 2:1) nécessite des stratégies de rééquilibrage pour éviter que le modèle
ne soit biaisé vers la classe majoritaire. La technique \textbf{SMOTE} (\textit{Synthetic Minority
Over-sampling Technique}) combinée à \texttt{class\_weight='balanced'} a été retenue.

\subsection{Analyse des Corrélations et Fuite de Données}

L'analyse exploratoire a révélé que plusieurs features présentaient des corrélations anormalement
élevées avec la variable cible, indiquant une \textbf{fuite de données} (voir Section~\ref{sec:leakage}).

% ==============================================================================
\section{Prétraitement des Données}
% ==============================================================================

\subsection{Pipeline de Nettoyage (6 Étapes)}

Le module \texttt{preprocessing.py} applique un pipeline structuré en 6 étapes séquentielles :

\begin{enumerate}[label=\textbf{Étape \arabic*.}]

  \item \textbf{Suppression des features à variance nulle} \\
  Élimination automatique des colonnes constantes qui n'apportent aucune information discriminante.

  \item \textbf{Correction des valeurs aberrantes codées} \\
  Remplacement par \texttt{NaN} des codes invalides identifiés dans les données sources :
  \begin{itemize}
    \item \texttt{SupportTicketsCount} : valeurs $-1$ et $999$
    \item \texttt{SatisfactionScore} : valeurs $-1$ et $99$
    \item \texttt{Age} : valeurs $-1$ et $999$
  \end{itemize}

  \item \textbf{Traitement des dates d'inscription} \\
  La feature \texttt{RegistrationDate} est décomposée en 4 features temporelles :
  ancienneté en jours, mois d'inscription, trimestre, indicateur de week-end.

  \item \textbf{Feature engineering sur les adresses IP} \\
  Extraction de features réseau à partir de \texttt{LastLoginIP} (classe réseau, plage d'adresses).

  \item \textbf{Création de features comportementales dérivées} \\
  Calcul de \texttt{AvgBasketValue}, \texttt{ReturnRatio}, \texttt{CancelRate},
  \texttt{AvgDaysBetweenPurchases}, \texttt{WeekendPurchaseRatio}, etc.

  \item \textbf{Encodage des variables catégorielles} \\
  Trois stratégies combinées :
  \begin{itemize}
    \item \textbf{Encodage ordinal} : \texttt{AgeCategory}, \texttt{LoyaltyLevel},
          \texttt{ChurnRiskCategory}, \texttt{BasketSizeCategory}, \texttt{PreferredTimeOfDay}
    \item \textbf{One-Hot Encoding} : \texttt{RFMSegment}, \texttt{CustomerType},
          \texttt{FavoriteSeason}, \texttt{Region}, \texttt{Gender}, etc.
    \item \textbf{Encodage binaire} : \texttt{NewsletterSubscribed}
  \end{itemize}
\end{enumerate}

Le résultat est un DataFrame de \textbf{87 features} (contre 52 initiales) après encodage.

\subsection{Gestion des Valeurs Manquantes}

L'imputation est intégrée \textbf{directement dans les pipelines scikit-learn} via
\texttt{KNNImputer(n\_neighbors=5)}, ce qui garantit que les statistiques d'imputation
sont calculées uniquement sur les données d'entraînement et appliquées aux données de test
— évitant ainsi une fuite d'information lors de l'imputation.

\subsection{Gestion du Déséquilibre de Classes}

\begin{table}[H]
\centering
\caption{Stratégies adoptées pour le déséquilibre}
\begin{tabular}{ll}
\toprule
\textbf{Technique} & \textbf{Paramètre} \\
\midrule
SMOTE & \texttt{random\_state=42} \\
Poids de classe & \texttt{class\_weight='balanced'} \\
Seuil de décision & Abaissé de 0.5 → 0.35 (recommandé) \\
\bottomrule
\end{tabular}
\end{table}

% ==============================================================================
\section{Problème de Data Leakage}
\label{sec:leakage}
% ==============================================================================

\subsection{Symptôme Détecté}

Lors des premières expériences, le modèle de classification atteignait un AUC-ROC de \textbf{1.00} —
une perfection irréaliste qui signale inévitablement une \textbf{fuite de données} (\textit{data leakage}).

\begin{warnbox}[Data Leakage Détecté]
AUC initial = \textbf{1.00} $\rightarrow$ Performance irréaliste et non généralisable.\\
Cause : présence de features contenant des informations du futur ou directement dérivées de la cible.
\end{warnbox}

\subsection{Features Fuyantes Identifiées}

L'analyse des corrélations a révélé 14+ features problématiques :

\begin{table}[H]
\centering
\caption{Features responsables de la fuite de données}
\begin{tabular}{lll}
\toprule
\textbf{Feature} & \textbf{Corrélation avec Churn} & \textbf{Raison} \\
\midrule
\texttt{ChurnRiskCategory} & 0.88 & Construite pour prédire le churn \\
\texttt{Recency} & 0.86 & Client parti $\Rightarrow$ récence élevée \\
\texttt{CustomerType\_Perdu} & 0.70 & « Perdu » = churné par définition \\
\texttt{RFMSegment\_Dormants} & 0.58 & « Dormant » = inactif \\
\texttt{LoyaltyLevel} & $-$0.43 & Calculé sur historique récent \\
\texttt{AccountStatus\_Closed} & — & Statut post-churn \\
\texttt{TenureRatio} & — & Dérivé de la date de résiliation \\
\bottomrule
\end{tabular}
\end{table}

\subsection{Solution Appliquée}

Exclusion stricte de toutes les features fuyantes via la liste \texttt{LEAKY\_CHURN\_FEATURES}
dans \texttt{preprocessing.py} :

\begin{lstlisting}[language=Python, caption=Exclusion des features fuyantes]
LEAKY_CHURN_FEATURES = [
    'Recency', 'CustomerTenureDays', 'FirstPurchaseDaysAgo',
    'TenureRatio', 'MonetaryPerDay',
    'ChurnRiskCategory', 'LoyaltyLevel', 'SpendingCategory',
    'CustomerType_Perdu', 'CustomerType_Hyperactif',
    'CustomerType_Nouveau', 'CustomerType_Occasionnel',
    'CustomerType_Regulier',
    'RFMSegment_Champions', 'RFMSegment_Dormants',
    'RFMSegment_Fideles', 'RFMSegment_Potentiels',
    'AccountStatus_Closed', 'AccountStatus_Suspended',
    'AccountStatus_Active', 'AccountStatus_Pending'
]
\end{lstlisting}

\begin{successbox}[Résultat après correction]
AUC-ROC : \textbf{1.00} $\rightarrow$ \textbf{0.958} \\
Le modèle délivre désormais des performances réalistes et généralisables à des données inédites.
\end{successbox}

% ==============================================================================
\section{Modèle 1 — Classification du Churn}
% ==============================================================================

\subsection{Objectif}

Détecter les clients à risque de départ \textbf{avant qu'ils ne partent}, afin de déclencher
des actions de rétention ciblées et précoces.

\subsection{Architecture du Pipeline}

Le pipeline intègre toutes les étapes de manière cohérente pour éviter toute fuite :

\begin{lstlisting}[language=Python, caption=Pipeline de classification (scikit-learn)]
ImbPipeline([
    ('imputer',    KNNImputer(n_neighbors=5)),
    ('scaler',     RobustScaler()),
    ('smote',      SMOTE(random_state=42)),
    ('classifier', RandomForestClassifier(
        n_estimators=200,
        max_depth=10,
        min_samples_leaf=10,
        class_weight='balanced'
    ))
])
\end{lstlisting}

Le choix de \texttt{RobustScaler} (plutôt que \texttt{StandardScaler}) assure une robustesse
aux valeurs aberrantes résiduelles. L'utilisation d'\texttt{ImbPipeline} (bibliothèque
\textit{imbalanced-learn}) permet d'intégrer SMOTE directement dans le pipeline,
garantissant que le suréchantillonnage n'est appliqué qu'aux données d'entraînement.

\subsection{Résultats et Métriques}

\begin{table}[H]
\centering
\caption{Performances du modèle de classification churn}
\begin{tabular}{lcc}
\toprule
\textbf{Métrique} & \textbf{Valeur} & \textbf{Interprétation} \\
\midrule
AUC-ROC      & \textbf{0.958} & Excellent pouvoir discriminant \\
Accuracy     & 88\%           & Précision globale \\
Precision (Churn=1) & 0.81  & 81\% des alertes sont réelles \\
Recall (Churn=1)    & \textbf{0.85} & 85\% des churners détectés \\
F1-Score     & 0.87           & Équilibre précision/rappel \\
\bottomrule
\end{tabular}
\end{table}

\begin{resultbox}[Interprétation Métier]
Le modèle identifie correctement \textbf{85\% des clients qui vont partir}.
En abaissant le seuil de décision de $0.5 \to 0.35$, il est possible d'augmenter ce rappel
au détriment de la précision — un arbitrage pertinent dans un contexte de rétention client
où manquer un churner coûte plus cher qu'une fausse alerte.
\end{resultbox}

\subsection{Importance des Features}

L'analyse des importances du Random Forest révèle les features les plus déterminantes :

\begin{table}[H]
\centering
\caption{Top 10 features d'importance pour la prédiction du churn}
\begin{tabular}{clc}
\toprule
\textbf{Rang} & \textbf{Feature} & \textbf{Importance (\%)} \\
\midrule
1 & \texttt{SatisfactionScore}          & 23.01\% \\
2 & \texttt{PreferredMonth}             & 17.25\% \\
3 & \texttt{Frequency}                  & 8.33\% \\
4 & \texttt{NegativeQuantityCount}      & 7.54\% \\
5 & \texttt{UniqueCountries}            & 4.89\% \\
6 & \texttt{WeekendPurchaseRatio}       & 4.31\% \\
7 & \texttt{TotalQuantity}              & 4.09\% \\
8 & \texttt{MonetaryTotal}              & 3.68\% \\
9 & \texttt{AgeCategory}               & 3.19\% \\
10 & \texttt{PreferredTimeOfDay}        & 3.00\% \\
\bottomrule
\end{tabular}
\end{table}

\textbf{Insight clé} : le score de satisfaction est la feature la plus prédictive (23\%),
ce qui souligne l'importance d'un suivi proactif de l'expérience client.

% ==============================================================================
\section{Modèle 2 — Segmentation Clients (KMeans)}
% ==============================================================================

\subsection{Objectif}

Regrouper les 4\,372 clients en \textbf{segments homogènes} basés sur leur comportement
RFM (\textit{Recency}, \textit{Frequency}, \textit{Monetary}) enrichi, afin de personnaliser
les actions marketing par profil.

\subsection{Prétraitement Spécifique au Clustering}

Le clustering nécessite un prétraitement adapté, distinct de celui de la classification :

\begin{enumerate}
  \item \textbf{Sélection de 11 features RFM} : \textit{Recency}, \textit{Frequency},
        \textit{MonetaryTotal}, \textit{SatisfactionScore}, \textit{ReturnRatio},
        \textit{UniqueProducts}, \textit{AvgDaysBetweenPurchases}, \textit{Age},
        \textit{SupportTicketsCount}, \textit{AvgBasketValue}, \textit{CancelRate}.
  \item \textbf{Filtrage} : exclusion des clients avec $\texttt{MonetaryTotal} = 0$
        (52 clients supprimés, 4\,320 retenus).
  \item \textbf{Winsorisation} aux percentiles 1--99 pour limiter l'influence des extrêmes.
  \item \textbf{Transformation logarithmique} (\texttt{np.log1p}) sur les variables asymétriques :
        \textit{Recency}, \textit{Frequency}, \textit{MonetaryTotal}, \textit{UniqueProducts},
        \textit{AvgDaysBetweenPurchases}, \textit{AvgBasketValue}.
  \item \textbf{Normalisation} avec \texttt{RobustScaler}.
  \item \textbf{Réduction de dimension} par ACP (20 composantes, expliquant la variance maximale).
\end{enumerate}

\subsection{Sélection du Nombre de Clusters}

La méthode du coude (\textit{elbow method}) combinée à l'analyse du score de silhouette a
orienté le choix vers $k = 4$ clusters, offrant le meilleur compromis entre cohésion intra-cluster
et séparation inter-cluster.

\subsection{Profils des 4 Segments}

\begin{table}[H]
\centering
\caption{Profils détaillés des 4 segments clients}
\begin{tabular}{lcccccc}
\toprule
\textbf{Segment} & \textbf{N} & \textbf{Monetary (£)} & \textbf{Recency (j)} & \textbf{Freq.} & \textbf{Churn} \\
\midrule
\textbf{VIP} ��         & 1\,311 & 5\,047 & 29 & 11.3 & 6.0\% \\
\textbf{Occasionnel}     & 1\,202 & 762   & 71 & 3.5  & 28.5\% \\
\textbf{Fidèle}          & 238   & 870   & 109 & 5.0 & 40.8\% \\
\textbf{À Risque} ⚠️    & 1\,569 & 365   & 154 & 1.2 & 57.1\% \\
\bottomrule
\end{tabular}
\end{table}

\subsubsection{Description des Segments}

\begin{itemize}
  \item \textbf{VIP (30\% de la base)} : clients à haute valeur, fréquents, récents.
        Taux de churn très bas (6\%). Ce sont les ambassadeurs de la marque.
  \item \textbf{Occasionnels (27.5\%)} : acheteurs à fréquence modérée,
        montant faible, churn modéré (28.5\%).
  \item \textbf{Fidèles (5.5\%)} : clients avec historique d'achat mais inactivité
        récente croissante — churn préoccupant (40.8\%).
  \item \textbf{À Risque (35.9\%)} : clients inactifs, faible dépense, forte
        probabilité de churn (57.1\%). Priorité d'action.
\end{itemize}

\subsection{Métriques de Qualité du Clustering}

\begin{table}[H]
\centering
\caption{Métriques de qualité du clustering KMeans}
\begin{tabular}{lll}
\toprule
\textbf{Métrique} & \textbf{Valeur} & \textbf{Interprétation} \\
\midrule
Silhouette Score   & 0.228   & Acceptable (données comportementales) \\
Davies-Bouldin     & 1.30    & Séparation raisonnable \\
Calinski-Harabasz  & 692     & Bonne cohésion intra-cluster \\
Inertie            & 20\,586\,522 & Compacité globale \\
\bottomrule
\end{tabular}
\end{table}

\begin{resultbox}[Note sur le Silhouette Score]
Un score de silhouette de 0.228 peut paraître faible sur des données comportementales continues
avec chevauchement naturel. La comparaison avec le score initial de 0.99 (obtenu sans
winsorisation ni log-transform) illustre qu'une valeur quasi-parfaite sur données non traitées
signalait des clusters triviaux dominés par les outliers — non de vrais groupes homogènes.
\end{resultbox}

% ==============================================================================
\section{Modèle 3 — Prévision de la Valeur Client}
% ==============================================================================

\subsection{Objectif}

Estimer le \textbf{montant total dépensé} (\texttt{MonetaryTotal}) par un client, afin de
scorer la base client par potentiel de revenus et d'allouer les ressources marketing en conséquence.

\subsection{Architecture du Pipeline}

\begin{lstlisting}[language=Python, caption=Pipeline de régression]
Pipeline([
    ('imputer',   KNNImputer(n_neighbors=5)),
    ('scaler',    RobustScaler()),
    ('regressor', LinearRegression())
])
\end{lstlisting}

Plusieurs modèles ont été évalués (Régression Linéaire, Ridge, Random Forest Regressor).
La Régression Linéaire a été retenue comme modèle de référence pour sa simplicité
et son interprétabilité, avec $R^2 = 0.668$.

\subsection{Résultats}

\begin{table}[H]
\centering
\caption{Performances du modèle de régression}
\begin{tabular}{lcc}
\toprule
\textbf{Métrique} & \textbf{Valeur} & \textbf{Interprétation} \\
\midrule
$R^2$  & \textbf{0.668} & 66.8\% de la variance expliquée \\
MAE    & 864 £          & Erreur absolue moyenne \\
RMSE   & 1\,124 £       & Erreur quadratique moyenne \\
\bottomrule
\end{tabular}
\end{table}

\begin{resultbox}[Exemple de Prédictions]
\begin{itemize}
  \item Client haute fréquence (80 commandes, 2\,999 £ dépensés) :\\
        Churn = \textbf{15.9\%}, Segment = \textit{VIP}, Valeur estimée = \textbf{88\,789 £}
  \item Client basse fréquence (faibles dépenses) :\\
        Churn = \textbf{96.4\%}, Segment = \textit{Occasionnel}, Valeur estimée = \textbf{3\,552 £}
\end{itemize}
\end{resultbox}

% ==============================================================================
\section{Application Flask — Déploiement}
% ==============================================================================

\subsection{Architecture de l'Application}

L'application Flask expose \textbf{4 pages principales} et une \textbf{API REST} :

\begin{table}[H]
\centering
\caption{Routes de l'application Flask}
\begin{tabular}{lll}
\toprule
\textbf{Route} & \textbf{Méthode} & \textbf{Description} \\
\midrule
\texttt{/}           & GET      & Dashboard — métriques et aperçu des segments \\
\texttt{/predict}    & GET/POST & Formulaire RFM → prédiction (churn + segment + valeur) \\
\texttt{/segments}   & GET      & Profils détaillés des 4 segments \\
\texttt{/about}      & GET      & Documentation technique et stack \\
\texttt{/api/predict}& POST     & Endpoint REST JSON \\
\texttt{/run-models} & POST     & Relance l'entraînement \\
\bottomrule
\end{tabular}
\end{table}

\subsection{API REST}

\subsubsection{Endpoint \texttt{POST /api/predict}}

\textbf{Requête JSON :}
\begin{lstlisting}[language=Python, caption=Exemple de requête API]
{
  "Frequency": 5,
  "MonetaryTotal": 1200,
  "Age": 35,
  "Gender": "M",
  "AgeCategory": "35-44",
  "PreferredMonth": 3,
  "FavoriteSeason": "Printemps",
  "SatisfactionScore": 4,
  "SupportTicketsCount": 1
}
\end{lstlisting}

\textbf{Réponse JSON :}
\begin{lstlisting}[language=Python, caption=Exemple de réponse API]
{
  "churn": 0,
  "probability": 0.15,
  "risk": "Faible",
  "cluster": 1,
  "cluster_name": "Fidele",
  "monetary_prediction": 1850.25
}
\end{lstlisting}

\subsection{Artefacts Générés}

L'entraînement produit 7 artefacts sauvegardés avec \texttt{joblib} :

\begin{itemize}
  \item \texttt{churn\_pipeline.pkl} — Pipeline complet Random Forest + SMOTE
  \item \texttt{clustering\_preprocessor\_pipeline.pkl} — KNNImputer + Scaler + PCA
  \item \texttt{kmeans\_model.pkl} — Modèle KMeans $k=4$
  \item \texttt{regression\_pipeline.pkl} — Pipeline Régression Linéaire
  \item \texttt{cluster\_label\_mapping.pkl} — Correspondance id $\to$ label métier
  \item \texttt{processed\_columns.pkl} — Colonnes attendues après preprocessing
  \item \texttt{dashboard\_metrics.pkl} — Métriques pré-calculées pour le dashboard
\end{itemize}

% ==============================================================================
\section{Résumé des Défis Résolus}
% ==============================================================================

\begin{table}[H]
\centering
\caption{Problèmes rencontrés et solutions apportées}
\begin{tabular}{p{3.5cm}p{2.5cm}p{2.5cm}p{4cm}}
\toprule
\textbf{Problème} & \textbf{Avant} & \textbf{Après} & \textbf{Solution} \\
\midrule
Data leakage churn & AUC = 1.00 & AUC = 0.958 & Exclusion de 14+ features fuyantes \\
\addlinespace
Clustering trivial & Silhouette = 0.99 & Silhouette = 0.228 & Winsorisation + log-transform \\
\addlinespace
Déséquilibre classes & 67/33\% & Équilibré & SMOTE + \texttt{class\_weight='balanced'} \\
\addlinespace
Fuite dans imputation & Avant split & Dans pipeline & KNNImputer intégré au Pipeline \\
\bottomrule
\end{tabular}
\end{table}

% ==============================================================================
\section{Stack Technique}
% ==============================================================================

\begin{table}[H]
\centering
\caption{Technologies utilisées dans le projet}
\begin{tabular}{ll}
\toprule
\textbf{Catégorie} & \textbf{Technologies} \\
\midrule
ML / Data Science & scikit-learn 1.3+, imbalanced-learn, pandas 2.0+, numpy 1.24+ \\
Visualisation & matplotlib, seaborn, plotly, Chart.js \\
Déploiement & Flask 2.3+, joblib, gunicorn \\
Frontend & Bootstrap 5.3, Font Awesome 6.4 \\
Développement & Jupyter Notebook (5 notebooks), pytest, black \\
Langage & Python 3.13 \\
\bottomrule
\end{tabular}
\end{table}

% ==============================================================================
\section{Recommandations Métier}
% ==============================================================================

\subsection{Réduction du Taux de Churn}

\begin{itemize}
  \item \textbf{Satisfaction client en priorité} : le score de satisfaction est la feature
        la plus prédictive (23\%). Un programme de suivi post-achat et de résolution rapide
        des tickets support peut réduire significativement le churn.
  \item \textbf{Campagnes saisonnières ciblées} : \texttt{PreferredMonth} est la 2\textsuperscript{e}
        feature la plus importante. Des offres personnalisées selon la saisonnalité préférée
        du client sont recommandées.
  \item \textbf{Ajustement du seuil de décision} : passer de $0.5 \to 0.35$ pour maximiser
        la détection des churners au prix d'une légère baisse de précision — arbitrage favorable
        quand le coût de rétention est inférieur au coût d'acquisition d'un nouveau client.
\end{itemize}

\subsection{Personnalisation Marketing par Segment}

\begin{itemize}
  \item \textbf{VIP (churn 6\%)} : programme premium, ventes privées, accès anticipé
        aux nouvelles collections, service client dédié.
  \item \textbf{Fidèles (churn 40.8\%)} : campagnes de réactivation, cross-selling ciblé,
        programme de fidélité avec récompenses progressives.
  \item \textbf{Occasionnels (churn 28.5\%)} : promotions saisonnières, offres de bienvenue,
        newsletters personnalisées pour augmenter la fréquence d'achat.
  \item \textbf{À Risque (churn 57.1\%)} : enquête de satisfaction, geste commercial,
        offre de réengagement urgente avant la résiliation définitive.
\end{itemize}

\subsection{Optimisation du Chiffre d'Affaires}

\begin{itemize}
  \item Scorer les clients par potentiel de dépense avant chaque campagne marketing.
  \item Envisager un modèle XGBoost dédié aux clients à dépenses élevées (outliers).
  \item Collecter des données de navigation et de wishlist pour enrichir les features
        comportementales et améliorer le $R^2$ du modèle de régression.
\end{itemize}

% ==============================================================================
\section{Améliorations Futures}
% ==============================================================================

\subsection{Court Terme}

\begin{itemize}
  \item Hyperparameter tuning via GridSearchCV ou Optuna (notamment pour Random Forest et KMeans).
  \item Feature selection automatique (RFE, SHAP values) pour réduire la dimensionnalité.
  \item Tests unitaires avec pytest pour garantir la robustesse du pipeline.
  \item Documentation API avec Swagger / OpenAPI.
\end{itemize}

\subsection{Moyen Terme}

\begin{itemize}
  \item Déploiement Docker + Kubernetes pour la mise en production scalable.
  \item Monitoring du data drift avec Evidently AI (détection des dégradations de performance).
  \item Dashboard interactif avec Streamlit ou Dash pour les équipes métier.
  \item A/B testing des modèles en production pour valider l'impact réel sur la rétention.
\end{itemize}

\subsection{Long Terme}

\begin{itemize}
  \item Deep Learning pour la modélisation des séquences temporelles d'achats (LSTM, Transformer).
  \item Pipeline MLOps complet avec MLflow (tracking des expériences) et DVC (versioning des données).
  \item Prédictions en temps réel via Apache Kafka et streaming.
  \item AutoML (H2O.ai) pour l'exploration automatisée des architectures de modèles.
\end{itemize}

% ==============================================================================
\section{Conclusion}
% ==============================================================================

Ce projet démontre la mise en œuvre d'un \textbf{pipeline Machine Learning complet et
production-ready} pour l'analyse comportementale d'une clientèle e-commerce. Les trois
objectifs principaux ont été atteints :

\begin{enumerate}
  \item \textbf{Classification du churn} avec un AUC-ROC de \textbf{0.958}, après correction
        d'un problème de data leakage qui produisait initialement un score parfait et trompeur de 1.00.
  \item \textbf{Segmentation en 4 clusters} (VIP, Fidèle, Occasionnel, À Risque) avec des
        profils métier distincts et actionnables, suite à un prétraitement rigoureux par
        winsorisation et transformation logarithmique.
  \item \textbf{Prévision de la valeur client} avec un $R^2$ de \textbf{0.668}, permettant
        de scorer la base client par potentiel de revenus.
\end{enumerate}

L'ensemble du système est encapsulé dans une \textbf{application Flask} avec interface graphique
et API REST, rendant les prédictions accessibles aux équipes métier sans expertise technique.

Les défis résolus — notamment la détection et correction du data leakage, la gestion du
déséquilibre de classes par SMOTE, et le clustering non trivial par transformation logarithmique —
illustrent les pièges classiques du Machine Learning en production et les bonnes pratiques
pour les surmonter.

\bigskip
\begin{center}
{\color{primary}\rule{0.6\linewidth}{0.8pt}}\\[0.4em]
{\small\itshape
Rapport réalisé dans le cadre du Module Machine Learning — GI2S1 MolkaChaoui — 2025--2026\\
Atelier : Analyse Comportementale Clientèle Retail}
\end{center}

\end{document}